% answers_short.tex —— 1994 年·上海卷·数学 · 简短答案速查
% 由 tools/make_short_answers.py 自 answers.tex 的 \gksolhead/\gkanswer 提取，
% 并与 1994/1994_数学真题_上海_答案要点.md 逐题比对后生成，勿手改（改源 MD 后重跑）。

\begin{gknotice}
\textbf{使用说明}：本节只给最终答案便于快速核对，完整推导见解析版。本卷第1—20题为项目整理；第21—24、26题据固定来源所附解答整理。第25题项目推导经独立盲核；固定来源原印载的$t\ge2$分段仍保留，非官方答案。整卷其余内容尚未经过异人盲核。\textbf{非官方标准答案}。
\end{gknotice}

\gkpart{填空题}
\begin{itemize}[leftmargin=2.6em]
  \item[\textbf{1.}] $-2<x<0$
  \item[\textbf{2.}] $\dfrac{\sqrt7}{4}$
  \item[\textbf{3.}] $(x+2)^2+(y-3)^2=4$
  \item[\textbf{4.}] $x=4$
  \item[\textbf{5.}] $24$ 元
  \item[\textbf{6.}] $y=-\sqrt{x^2+1}\ (x\ge0)$
  \item[\textbf{7.}] $(0,\sqrt3)$ 和 $(0,-\sqrt3)$
  \item[\textbf{8.}] $40$
  \item[\textbf{9.}] $\sqrt2-1$
  \item[\textbf{10.}] $\dfrac\pi2$
\end{itemize}

\gkpart{选择题}
\begin{center}\begin{tabular}{ccccccccccc}
\toprule
题号 & 11 & 12 & 13 & 14 & 15 & 16 & 17 & 18 & 19 & 20 \\
\midrule
答案 & B & A & C & C & D & B & A & D & B & C \\
\bottomrule
\end{tabular}\end{center}

\gkpart{解答题}
\begin{itemize}[leftmargin=2.6em]
  \item[\textbf{21.}] $\dfrac7{24}$。
  \item[\textbf{22.}] $w=-\sqrt2+\sqrt2i$。
  \item[\textbf{23.}] (1) 二面角 $P-CD-A$ 的大小为 $\arctan\dfrac{\sqrt5}{3}$；(2) 点 $A$ 到平面 $PBC$ 的距离为 $\dfrac{\sqrt2}{2}a$
  \item[\textbf{24.}] (1) $t^2(t^2-1)x^2+(t^2-1)y^2=t^2$；(2) 轨迹为抛物线 $x^2=\dfrac{\sqrt2}{2}y$ 在直线 $x=\dfrac{\sqrt2}{2}$ 右侧的部分，以及抛物线 $x^2=-\dfrac{\sqrt2}{2}y$ 在直线 $x=-\dfrac{\sqrt2}{2}$ 左侧的部分
  \item[\textbf{25.}] (1) $S(t)=\begin{cases}2(1-t+t^2-t^3),&0<t<\dfrac12,\\[3pt]\dfrac12\left(t+\dfrac1t\right),&t\ge\dfrac12.\end{cases}$；(2) $S(t)$ 在 $(0,1]$ 上递减、在 $[1,+\infty)$ 上递增，最小值为 $S(1)=1$
  \item[\textbf{26.}] (1) $0<q<\dfrac{1+\sqrt5}{2}$；(2) $b_n=(1+r)q^{n-1}$，且 $\displaystyle\lim_{n\to\infty}\frac1{S_n}=\begin{cases}\dfrac{1-q}{1+r},&0<q<1,\\[3pt]0,&q\ge1.\end{cases}$；(3) 最大项为 $\dfrac94$，最小项为 $-4$
\end{itemize}
